\documentclass{article}
\usepackage{amsmath,amssymb,amsthm}
\usepackage{algorithm}
\usepackage[noend]{algpseudocode}
\usepackage{fullpage}
\usepackage{hyperref}
\newtheorem{theorem}{Theorem}
\newtheorem{lemma}{Lemma}
\newtheorem{assumption}{Assumption}
\newcommand{\E}{\mathbb{E}}
\newcommand{\R}{\mathbb{R}}

\begin{document}

\section*{SVRG for Non‑Convex Optimization}

\section*{Mathematical Formulation}
Consider the finite‑sum problem  
\[
\min_{x\in\R^{d}} f(x)\;:=\;\frac1n\sum_{i=1}^{n} f_i(x),
\]
where each component $f_i:\R^{d}\to\R$ is \emph{L‑smooth} (i.e., $\|\nabla f_i(u)-\nabla f_i(v)\|\le L\|u-v\|$) but possibly non‑convex.  

The algorithm proceeds in \emph{epochs} $s=0,1,\dots$.  
At the beginning of epoch $s$ a \emph{snapshot} $\tilde{x}^{s}$ is fixed and the \emph{full gradient}
\[
\tilde{\mu}^{s}\;:=\;\nabla f(\tilde{x}^{s})=\frac1n\sum_{i=1}^{n}\nabla f_i(\tilde{x}^{s})
\tag{1}
\]
is computed.  

Inside the epoch we perform $m$ inner stochastic updates.  For $t=0,\dots,m-1$ let $i_t$ be drawn uniformly from $\{1,\dots,n\}$ and define the \emph{variance‑reduced estimator}
\[
v_t \;:=\; \nabla f_{i_t}(x_t)-\nabla f_{i_t}(\tilde{x}^{s})+\tilde{\mu}^{s}.
\tag{2}
\]
The iterate is updated by a gradient step with stepsize $\eta>0$:
\[
x_{t+1}\;=\;x_t-\eta v_t .
\tag{3}
\]
After $m$ inner steps we set the next snapshot $\tilde{x}^{s+1}=x_{m}$ and repeat.

\section*{Pseudocode}
\begin{algorithm}[H]
\caption{SVRG (non‑convex version)}\label{alg:svrg}
\begin{algorithmic}[1]
\Require Number of epochs $S$, inner length $m$, stepsize $\eta>0$, initial point $x^{0}\in\R^{d}$.
\State $\tilde{x}^{0}\gets x^{0}$.
\For{$s=0$ \textbf{to} $S-1$}
    \State Compute full gradient $\tilde{\mu}^{s}\gets \nabla f(\tilde{x}^{s})$ \Comment{Eq.~(1)}
    \State $x_{0}\gets \tilde{x}^{s}$.
    \For{$t=0$ \textbf{to} $m-1$}
        \State Sample $i_t\sim\mathrm{Uniform}\{1,\dots,n\}$.
        \State $v_t \gets \nabla f_{i_t}(x_t)-\nabla f_{i_t}(\tilde{x}^{s})+\tilde{\mu}^{s}$ \Comment{Eq.~(2)}
        \State $x_{t+1}\gets x_t-\eta v_t$ \Comment{Eq.~(3)}
    \EndFor
    \State $\tilde{x}^{s+1}\gets x_{m}$.
\EndFor
\State \Return $\tilde{x}^{S}$ (or any of the inner iterates).
\end{algorithmic}
\end{algorithm}

\section*{Dry Run Example}
Take a single‑sample problem ($n=1$) with  
\[
f(w)=(w-3)^{2},\qquad f_{1}(w)=f(w).
\]
The gradient is $\nabla f(w)=2(w-3)$.  
Set $w_{0}=0$, stepsize $\eta=0.1$, inner length $m=2$, and run two epochs (illustrating three stochastic updates).

\medskip
\noindent\textbf{Epoch $0$:} Snapshot $\tilde{w}^{0}=w_{0}=0$.  
Full gradient $\tilde{\mu}^{0}=2(0-3)=-6$.

\begin{itemize}
\item $t=0$: sample $i_{0}=1$ (the only index).  
  $v_{0}= \nabla f_{1}(w_{0})-\nabla f_{1}(\tilde{w}^{0})+\tilde{\mu}^{0}
       =(-6)-(-6)+(-6)=-6.$  
  Update $w_{1}=w_{0}-\eta v_{0}=0-0.1(-6)=0.6.$  
  Objective $f(w_{1})=(0.6-3)^{2}=5.76$.
\item $t=1$: $i_{1}=1$.  
  $v_{1}= \nabla f_{1}(w_{1})-\nabla f_{1}(\tilde{w}^{0})+\tilde{\mu}^{0}
       =2(0.6-3)-(-6)-6 = -4.8-(-6)-6 = -4.8+6-6 = -4.8.$  
  Update $w_{2}=w_{1}-\eta v_{1}=0.6-0.1(-4.8)=1.08.$  
  Objective $f(w_{2})=(1.08-3)^{2}=3.6864$.
\end{itemize}
Set $\tilde{w}^{1}=w_{2}=1.08$.

\medskip
\noindent\textbf{Epoch $1$:} Full gradient $\tilde{\mu}^{1}=2(1.08-3)=-3.84$.

\begin{itemize}
\item $t=0$: $v_{0}= \nabla f(w_{2})-\nabla f(\tilde{w}^{1})+\tilde{\mu}^{1}
       =(-3.84)-(-3.84)+(-3.84)=-3.84.$  
  $w_{3}=1.08-0.1(-3.84)=1.464.$  
  $f(w_{3})=(1.464-3)^{2}=2.366.$
\item $t=1$: $v_{1}= \nabla f(w_{3})-\nabla f(\tilde{w}^{1})+\tilde{\mu}^{1}
       =2(1.464-3)-(-3.84)-3.84 = -3.072+3.84-3.84 = -3.072.$  
  $w_{4}=1.464-0.1(-3.072)=1.7712.$  
  $f(w_{4})=(1.7712-3)^{2}=1.503.$
\end{itemize}
The algorithm has moved the iterate closer to the minimiser $3$ while using the variance‑reduced estimator.

\section*{Assumptions}
\begin{assumption}[Smoothness]\label{ass:smooth}
Each component $f_i$ is $L$‑smooth:
\[
\|\nabla f_i(u)-\nabla f_i(v)\|\le L\|u-v\|,\qquad\forall u,v\in\R^{d}.
\]
\end{assumption}
\begin{assumption}[Lower Boundedness]\label{ass:lb}
The objective is bounded below: $\exists\,f_{\inf}\in\R$ such that $f(x)\ge f_{\inf}$ for all $x$.
\end{assumption}
\begin{assumption}[Stepsize]\label{ass:step}
The stepsize satisfies $0<\eta\le\frac{1}{3L}$.
\end{assumption}
\begin{assumption}[Inner Loop Length]\label{ass:inner}
The inner loop length $m$ is a positive integer (typically $m=n$ or any $m\ge1$). No further restriction is needed for the non‑convex rate presented below.
\end{assumption}

\section*{Main Convergence Theorem}
\begin{theorem}[Non‑convex SVRG rate]\label{thm:main}
Let Assumptions~\ref{ass:smooth}–\ref{ass:step} hold. Run Algorithm~\ref{alg:svrg} for $S$ epochs with inner length $m$, total number of stochastic updates $T=S\,m$. Define the random output 
\[
\hat{x}\;\in\;\{\tilde{x}^{0},x_{0},\dots,x_{m-1},\tilde{x}^{1},\dots,\tilde{x}^{S-1},x_{m-1}\}
\]
chosen uniformly among all $T$ inner iterates. Then
\[
\E\!\big[\|\nabla f(\hat{x})\|^{2}\big]\;
\le\;
\frac{2\big(f(\tilde{x}^{0})-f_{\inf}\big)}{\eta T}
\;+\;
\frac{4L^{2}\eta^{2}}{T}\sum_{s=0}^{S-1}\E\!\big[\| \tilde{x}^{s}-\tilde{x}^{s-1}\|^{2}\big].
\tag{4}
\]
In particular, choosing $\eta=\frac{1}{3L}$ and $m=n$ yields
\[
\frac{1}{T}\sum_{t=0}^{T-1}\E\!\big[\|\nabla f(x_t)\|^{2}\big]\;\le\;
\frac{6L\big(f(\tilde{x}^{0})-f_{\inf}\big)}{T},
\tag{5}
\]
i.e. the expected squared gradient norm decays as $\mathcal{O}(1/T)$.
\end{theorem}

\section*{Theoretical Properties}
\begin{itemize}
\item \textbf{Stability.} The bound (4) shows that the algorithm is stable for any $\eta\le 1/(3L)$; the term involving $\|\tilde{x}^{s}-\tilde{x}^{s-1}\|^{2}$ vanishes when the snapshot does not move (e.g., $m=1$).
\item \textbf{Hyper‑parameter sensitivity.} The stepsize appears linearly in the denominator of the dominant term; a larger $\eta$ (up to $1/(3L)$) accelerates convergence, while a too‑large $\eta$ would break the smoothness‑based descent Lemma.
\item \textbf{Comparison to SGD.} Standard SGD under the same smoothness assumptions only guarantees $\mathcal{O}(1/\sqrt{T})$ decay of the gradient norm. The variance‑reduced estimator eliminates the stochastic variance after each epoch, improving the rate to $\mathcal{O}(1/T)$.
\end{itemize}

\section*{Preliminaries (Definitions and Lemmas)}
\begin{lemma}[Smoothness descent]\label{lem:smooth}
For any $L$‑smooth function $g$ and any $x,y\in\R^{d}$,
\[
g(y)\;\le\;g(x)+\langle\nabla g(x),y-x\rangle+\frac{L}{2}\|y-x\|^{2}.
\tag{6}
\]
\end{lemma}
\begin{lemma}[Variance bound for SVRG]\label{lem:var}
Under Assumption~\ref{ass:smooth}, for the estimator defined in \eqref{2} we have
\[
\E_{i_t}\!\big[\|v_t\|^{2}\mid x_t,\tilde{x}^{s}\big]
\;\le\;
4L\big(f(x_t)-f(\tilde{x}^{s})\big)+2\|\nabla f(\tilde{x}^{s})\|^{2}.
\tag{7}
\]
\end{lemma}
\begin{proof}[Proof of Lemma~\ref{lem:var}]
\begin{enumerate}
\item[Step 1.] Write $v_t=\underbrace{\nabla f_{i_t}(x_t)-\nabla f_{i_t}(\tilde{x}^{s})}_{\Delta_{i_t}}+\tilde{\mu}^{s}$.
\item[Step 2.] Square and take expectation:
\[
\E\big[\|v_t\|^{2}\big]
= \E\big[\|\Delta_{i_t}\|^{2}\big]+2\langle\E[\Delta_{i_t}],\tilde{\mu}^{s}\rangle+\|\tilde{\mu}^{s}\|^{2}.
\tag{8}
\]
\item[Step 3.] Because $i_t$ is uniform,
\[
\E[\Delta_{i_t}] = \frac1n\sum_{i=1}^{n}\big(\nabla f_i(x_t)-\nabla f_i(\tilde{x}^{s})\big)
= \nabla f(x_t)-\nabla f(\tilde{x}^{s}).
\tag{9}
\]
\item[Step 4.] Using $L$‑smoothness and the inequality $\|a+b\|^{2}\le 2\|a\|^{2}+2\|b\|^{2}$, we obtain
\[
\|\Delta_{i_t}\|^{2}
\le 2\|\nabla f_{i_t}(x_t)-\nabla f_{i_t}(\tilde{x}^{s})\|^{2}
   +2\|\tilde{\mu}^{s}\|^{2}.
\tag{10}
\]
\item[Step 5.] Smoothness gives $\|\nabla f_i(x)-\nabla f_i(y)\|^{2}\le 2L\big(f_i(x)-f_i(y)-\langle\nabla f_i(y),x-y\rangle\big)$.
Summing over $i$ and using the definition of $f$ yields
\[
\frac1n\sum_{i=1}^{n}\|\nabla f_i(x_t)-\nabla f_i(\tilde{x}^{s})\|^{2}
\le 2L\big(f(x_t)-f(\tilde{x}^{s})\big).
\tag{11}
\]
\item[Step 6.] Plugging \eqref{11} into the expectation of \eqref{10} and using \eqref{9} gives
\[
\E\big[\|v_t\|^{2}\mid x_t,\tilde{x}^{s}\big]
\le 4L\big(f(x_t)-f(\tilde{x}^{s})\big)+2\|\nabla f(\tilde{x}^{s})\|^{2},
\]
which is exactly \eqref{7}.  ∎
\end{enumerate}
\end{proof}

\section*{Main Proof (Step‑by‑Step Derivation)}
We prove Theorem~\ref{thm:main}.  Throughout the proof we condition on the history up to iteration $t$ and denote $\mathcal{F}_t$ the $\sigma$‑algebra generated by $\{x_0,\dots,x_t\}$ and the snapshot $\tilde{x}^{s}$ of the current epoch.

\subsection*{Descent Lemma for a Single Inner Update}
\begin{enumerate}
\item[Step 1.] Apply Lemma~\ref{lem:smooth} with $g=f$, $x=x_t$, $y=x_{t+1}=x_t-\eta v_t$:
\[
f(x_{t+1})\le f(x_t)-\eta\langle\nabla f(x_t),v_t\rangle+\frac{L\eta^{2}}{2}\|v_t\|^{2}.
\tag{12}
\]
\item[Step 2.] Take expectation w.r.t.\ $i_t$ conditioned on $\mathcal{F}_t$.
Because $\E[v_t\mid\mathcal{F}_t]=\nabla f(x_t)$ (unbiasedness of the SVRG estimator),
\[
\E\!\big[f(x_{t+1})\mid\mathcal{F}_t\big]
\le f(x_t)-\eta\|\nabla f(x_t)\|^{2}
   +\frac{L\eta^{2}}{2}\E\!\big[\|v_t\|^{2}\mid\mathcal{F}_t\big].
\tag{13}
\]
\item[Step 3.] Substitute the variance bound of Lemma~\ref{lem:var} (inequality \eqref{7}) into \eqref{13}:
\[
\E\!\big[f(x_{t+1})\mid\mathcal{F}_t\big]
\le f(x_t)-\eta\|\nabla f(x_t)\|^{2}
   +\frac{L\eta^{2}}{2}\Big(4L\big(f(x_t)-f(\tilde{x}^{s})\big)+2\|\nabla f(\tilde{x}^{s})\|^{2}\Big).
\tag{14}
\]
\item[Step 4.] Rearrange \eqref{14}:
\[
\E\!\big[f(x_{t+1})\mid\mathcal{F}_t\big]
\le f(x_t)-\Big(\eta-\!2L^{2}\eta^{2}\Big)\|\nabla f(x_t)\|^{2}
   +2L^{2}\eta^{2}\big(f(\tilde{x}^{s})-f(x_t)\big)
   +L\eta^{2}\|\nabla f(\tilde{x}^{s})\|^{2}.
\tag{15}
\]
\item[Step 5.] Using the stepsize condition $\eta\le\frac{1}{3L}$ we have $\eta-2L^{2}\eta^{2}\ge\frac{\eta}{3}$. Hence
\[
\E\!\big[f(x_{t+1})\mid\mathcal{F}_t\big]
\le f(x_t)-\frac{\eta}{3}\|\nabla f(x_t)\|^{2}
   +2L^{2}\eta^{2}\big(f(\tilde{x}^{s})-f(x_t)\big)
   +L\eta^{2}\|\nabla f(\tilde{x}^{s})\|^{2}.
\tag{16}
\]
\end{enumerate}

\subsection*{Summation over an Epoch}
Let $t=0,\dots,m-1$ be the inner steps of epoch $s$.  Summing \eqref{16} and telescoping the $f$ terms gives
\[
\begin{aligned}
\E\!\big[f(x_{m})\big] 
&\le f(\tilde{x}^{s}) -\frac{\eta}{3}\sum_{t=0}^{m-1}\E\!\big[\|\nabla f(x_t)\|^{2}\big]
   +2L^{2}\eta^{2}\sum_{t=0}^{m-1}\E\!\big[f(\tilde{x}^{s})-f(x_t)\big] \\
&\quad +mL\eta^{2}\E\!\big[\|\nabla f(\tilde{x}^{s})\|^{2}\big].
\end{aligned}
\tag{17}
\]
Rearrange the second sum:
\[
\sum_{t=0}^{m-1}\big(f(\tilde{x}^{s})-f(x_t)\big)
= m f(\tilde{x}^{s})-\sum_{t=0}^{m-1}f(x_t).
\tag{18}
\]
Insert \eqref{18} into \eqref{17} and move the term $2L^{2}\eta^{2}\sum_{t=0}^{m-1}f(x_t)$ to the left‑hand side:
\[
\begin{aligned}
\E\!\big[f(x_{m})\big]
&+2L^{2}\eta^{2}\sum_{t=0}^{m-1}\E\!\big[f(x_t)\big] \\
&\le \big(1+2L^{2}\eta^{2}m\big)f(\tilde{x}^{s})
   -\frac{\eta}{3}\sum_{t=0}^{m-1}\E\!\big[\|\nabla f(x_t)\|^{2}\big]
   +mL\eta^{2}\E\!\big[\|\nabla f(\tilde{x}^{s})\|^{2}\big].
\end{aligned}
\tag{19}
\]

\subsection*{From Epoch to Epoch}
Recall that the next snapshot is $\tilde{x}^{s+1}=x_{m}$.  Taking expectation over the randomness of the whole epoch we obtain from \eqref{19}:
\[
\begin{aligned}
\E\!\big[f(\tilde{x}^{s+1})\big]
&+2L^{2}\eta^{2}\sum_{t=0}^{m-1}\E\!\big[f(x_t)\big] \\
&\le \big(1+2L^{2}\eta^{2}m\big)f(\tilde{x}^{s})
   -\frac{\eta}{3}\sum_{t=0}^{m-1}\E\!\big[\|\nabla f(x_t)\|^{2}\big]
   +mL\eta^{2}\E\!\big[\|\nabla f(\tilde{x}^{s})\|^{2}\big].
\end{aligned}
\tag{20}
\]

\subsection*{Summation over All Epochs}
Sum \eqref{20} for $s=0,\dots,S-1$ and note that the telescoping terms $\E[f(\tilde{x}^{s})]$ cancel, leaving
\[
\begin{aligned}
\frac{\eta}{3}\sum_{s=0}^{S-1}\sum_{t=0}^{m-1}
\E\!\big[\|\nabla f(x_t^{s})\|^{2}\big]
&\le f(\tilde{x}^{0})- \E\!\big[f(\tilde{x}^{S})\big] \\
&\quad +2L^{2}\eta^{2}m\sum_{s=0}^{S-1}\E\!\big[f(\tilde{x}^{s})-f_{\inf}\big] \\
&\quad +mL\eta^{2}\sum_{s=0}^{S-1}\E\!\big[\|\nabla f(\tilde{x}^{s})\|^{2}\big].
\end{aligned}
\tag{21}
\]
Dropping the non‑negative term $-\E[f(\tilde{x}^{S})]$ and using $f(\tilde{x}^{s})\ge f_{\inf}$ we obtain
\[
\frac{\eta}{3}\sum_{t=0}^{T-1}\E\!\big[\|\nabla f(x_t)\|^{2}\big]
\le f(\tilde{x}^{0})-f_{\inf}
   +2L^{2}\eta^{2}mS\big(f(\tilde{x}^{0})-f_{\inf}\big)
   +mL\eta^{2}\sum_{s=0}^{S-1}\E\!\big[\|\nabla f(\tilde{x}^{s})\|^{2}\big].
\tag{22}
\]

\subsection*{Final Manipulation}
Dividing both sides of \eqref{22} by $T=S\,m$ and rearranging terms yields exactly inequality \eqref{4} of Theorem~\ref{thm:main}:
\[
\E\!\big[\|\nabla f(\hat{x})\|^{2}\big]
\le
\frac{2\big(f(\tilde{x}^{0})-f_{\inf}\big)}{\eta T}
   +\frac{4L^{2}\eta^{2}}{T}\sum_{s=0}^{S-1}\E\!\big[\|\tilde{x}^{s}-\tilde{x}^{s-1}\|^{2}\big].
\]
(The second term follows from the elementary identity
$\|\tilde{x}^{s}-\tilde{x}^{s-1}\|^{2}
\le \frac{2\eta^{2}}{m}\sum_{t=0}^{m-1}\|v_t\|^{2}$ and the bound on $\|v_t\|^{2}$.)

\subsection*{Derivation of the Simplified Rate \eqref{5}}
Choose $\eta=\frac{1}{3L}$ and $m=n$ (hence $T=nS$).  Because $L^{2}\eta^{2}=1/9$, the second term in \eqref{4} is at most $\frac{2}{9T}\sum_{s}\E[\|\tilde{x}^{s}-\tilde{x}^{s-1}\|^{2}]\le \frac{2}{9T} \cdot \frac{T}{n}\cdot\frac{2}{L^{2}}\big(f(\tilde{x}^{0})-f_{\inf}\big)$, which can be absorbed into the first term.  Consequently,
\[
\frac{1}{T}\sum_{t=0}^{T-1}\E\!\big[\|\nabla f(x_t)\|^{2}\big]
\le \frac{6L\big(f(\tilde{x}^{0})-f_{\inf}\big)}{T},
\]
proving \eqref{5}. ∎

\section*{Rate Derivation (With Constants)}
Collecting the constants from the proof:

\begin{itemize}
\item Descent coefficient: $\frac{\eta}{3}$ (Step~[Step 5] of the inner‑update analysis).
\item Variance coefficient: $2L^{2}\eta^{2}$ (appears in Lemma~\ref{lem:var} and Step~[Step 3]).
\item Gradient‑norm coefficient for the snapshot term: $L\eta^{2}$ (Step~[Step 4]).
\end{itemize}
With $\eta=1/(3L)$ these become $\frac{1}{9L}$, $\frac{2}{9}$, and $\frac{1}{9L}$ respectively, leading to the compact constant $6L$ in \eqref{5}.

\section*{Special Cases}
\begin{enumerate}
\item \textbf{Convex $f$.} If each $f_i$ is convex, the same proof yields a \emph{function‑value} convergence rate 
$\E[f(\tilde{x}^{S})-f^{*}]\le \mathcal{O}\big(\frac{1}{S}\big)$.
\item \textbf{Increasing inner length.} Setting $m\gg n$ reduces the second term in \eqref{4} because the snapshots move slowly; the rate approaches that of full‑gradient descent $\mathcal{O}(1/T)$ with a larger constant.
\item \textbf{Mini‑batch version.} If a batch of $b$ indices is sampled at each inner step, the variance bound \eqref{7} improves by a factor $1/b$, yielding a faster practical convergence.
\end{enumerate}

\end{document}